\subsection*{Questão 04}

\begin{questionbox}{04}
  Um processador executa tarefas que são compostas por uma sequência de micro-operações. Existem dois tipos de micro-operações: tipo A, que dura 1 ciclo de clock, e tipo B, que dura 2 ciclos de clock. Uma tarefa é definida pela sequência de micro-operações que a compõe.

  Quantas sequências distintas de micro-operações podem formar uma tarefa com duração total de exatamente 8 ciclos de clock?
\end{questionbox}

\begin{solutionbox}
  Seja $f(n)$ o número de sequências que totalizam $n$ ciclos de clock.

  Uma sequência que soma $n$ pode terminar em:
  \begin{itemize}
    \item uma operação do tipo A (1 ciclo), restando $n-1$ ciclos;
    \item uma operação do tipo B (2 ciclos), restando $n-2$ ciclos.
  \end{itemize}

  Logo, temos a recorrência:
  \[
  f(n) = f(n-1) + f(n-2).
  \]

  Com condições iniciais:
  \[
  f(0) = 1, \quad f(1) = 1.
  \]
  \end{solutionbox}

  \begin{solutionbox}
  Calculando iterativamente:
  \[
  f(2)=2,\; f(3)=3,\; f(4)=5,\; f(5)=8,\; f(6)=13,\; f(7)=21,\; f(8)=34.
  \]

  Portanto, o número de sequências distintas é:
  \[
  \boxed{34}.
  \]
\end{solutionbox}